\mypara{Primed Variables}
Another useful feature is for a variable to reference its own value from a previous context. We implemented this using ``primed''
variables. A primed variable is denoted as a regular variable followed by an apostrophe, and they can only be used as arguments for 
predicates.

\begin{figure}
    \begin{smallhskcode}
groupTest = test "group"
        (group :: [Int] -> [[Int]])
        [ex| ({a => ne a'} -> {b => gt <0>} -> a*{b})*{c}
            == [a*{b}]*{c} |]
    \end{smallhskcode}
    \caption{Primed variable in groupTest}
    \label{fig:group-primed}
\end{figure}

When evaluating a primed variable we consider the nearest scope. In the case of the the primed ``$a$'' in Figure~\ref{fig:group-primed},
the nearest scope is the label ``$c$'' on the kleene closure surrounding it. When $c=0$, we return an arbitrary value for $a'$,
for all $c > 0$, we evaluate $a'$ as the value stored for $a$ when $c$ equalled $n-1$. This directive can also be used to construct
a sorted list by making sure each element is less than the previous. 
% In a broad sense, it helps maintain a certain order which can be helpful when
% trying to write good tests for data structures that aren't concerned with order such as maps and sets. \evan{maybe don't even include that last bit...}

